% ----------------------------------------------------------
% Introdução
% ----------------------------------------------------------
\chapter{Introdução}
\label{cap:introducao}
% ----------------------------------------------------------

O setor automotivo possui uma quantidade significativa de informações técnicas registradas em manuais, catálogos, guias de uso e documentos de manutenção. Esses materiais reúnem orientações sobre fluidos, pneus, luzes de advertência, localização de componentes, procedimentos de verificação e cuidados de segurança. Apesar disso, o acesso prático a essas informações nem sempre é simples. Um usuário que deseja saber onde fica a tampa de abastecimento do óleo, qual especificação de óleo utilizar ou como interpretar um aviso no painel pode precisar consultar um documento em \acf{PDF} com centenas de páginas, procurar a seção correta e ainda interpretar termos técnicos utilizados pela montadora.

Essa dificuldade motivou o desenvolvimento deste trabalho. A proposta não é substituir o manual do veículo, mas construir uma forma mais direta de consultá-lo. Para isso, foi desenvolvido um sistema baseado em \acf{RAG}, no qual uma pergunta feita em linguagem natural é usada para recuperar trechos relevantes de manuais automotivos e, em seguida, gerar uma resposta fundamentada nesse contexto. Diferentemente de um chatbot que responde apenas a partir do conhecimento armazenado nos parâmetros de um modelo de linguagem, o sistema busca primeiro as evidências nos documentos disponíveis.

O título deste trabalho, \textit{Uso de Retrieval-Augmented Generation para suporte informacional no setor automotivo. Uma análise comparativa com grandes modelos de linguagem}, expressa duas dimensões do projeto. A primeira é a construção de um artefato computacional voltado ao suporte informacional, com uso de manuais oficiais e execução local. A segunda é a preparação de uma análise comparativa com grandes modelos de linguagem quando estes não possuem acesso aos mesmos manuais. Essa comparação é parte central da proposta, mas ainda não foi concluída nesta etapa, pois o sistema continua em processo de refinamento e validação.

\section{Problema de pesquisa}
\label{sec:problema}

O problema investigado pode ser descrito da seguinte forma: como utilizar recuperação documental e modelos de linguagem para tornar informações técnicas automotivas mais acessíveis, rastreáveis e contextualizadas para usuários que consultam manuais de veículos?

Esse problema envolve mais do que recuperar uma página do manual. Em uma interação conversacional, o sistema precisa compreender o veículo citado, manter o contexto entre mensagens, diferenciar uma pergunta nova de um complemento e evitar respostas que apenas encaminhem o usuário de volta ao manual. Por exemplo, se o usuário pergunta ``Onde fica a tampa do óleo do Audi A4?'' e depois informa apenas ``2008'', a segunda mensagem deve completar a intenção anterior. A busca documental deve receber uma pergunta reescrita, como ``Onde fica a tampa do óleo do Audi A4 2008?''.

Outro desafio é visual. Algumas perguntas exigem diagramas ou imagens, mas não basta retornar qualquer figura da mesma página. Uma imagem de abertura do capô, por exemplo, pode estar próxima de um trecho sobre o compartimento do motor, mas não responde onde fica a tampa de abastecimento do óleo. Por isso, o sistema precisa validar se a imagem realmente corresponde à peça ou ao procedimento perguntado.

\section{Objetivos}
\label{sec:objetivos}

O objetivo geral deste trabalho é desenvolver e documentar um sistema de suporte informacional automotivo baseado em RAG, capaz de consultar manuais, recuperar evidências textuais e visuais, gerar respostas fundamentadas e preparar uma avaliação comparativa com grandes modelos de linguagem sem acesso aos manuais.

\subsection{Objetivo geral}

Desenvolver uma aplicação web que permita consultar manuais automotivos por meio de perguntas em linguagem natural, utilizando recuperação documental, geração local de respostas, memória conversacional e seleção de imagens relevantes.

\subsection{Objetivos específicos}

\begin{itemize}
  \item organizar uma base local de manuais automotivos de diferentes modelos, anos e montadoras;
  \item implementar um pipeline de extração de texto, imagens e metadados a partir dos PDFs;
  \item construir chunks documentais preservando páginas, seções, entidades e imagens associadas;
  \item indexar os textos com embeddings e FAISS;
  \item combinar busca vetorial, BM25, regras de domínio e metadados do veículo;
  \item integrar o RAG a uma aplicação com backend NestJS e frontend Next.js;
  \item executar localmente a geração de respostas com Ollama e DeepSeek R1 8B;
  \item implementar memória conversacional e reescrita de consultas;
  \item selecionar imagens por meio de CLIP local e validação de relevância;
  \item propor uma metodologia de comparação entre o sistema desenvolvido e grandes modelos de linguagem sem acesso aos manuais.
\end{itemize}

\section{Justificativa}
\label{sec:justificativa}

A justificativa principal está na necessidade de aproximar o usuário das informações técnicas já existentes. Em muitos casos, a resposta correta está no manual, mas o formato do documento dificulta a consulta rápida. Um assistente conversacional com acesso aos manuais pode reduzir essa barreira, desde que mantenha rastreabilidade das fontes e não produza respostas inventadas.

O uso de RAG é adequado nesse cenário porque permite que o conhecimento do sistema seja atualizado pela base documental, e não apenas pelo treinamento do modelo. Isso é importante em domínio automotivo, pois as recomendações podem variar conforme ano, versão, motorização e mercado. Além disso, a execução local foi escolhida para manter maior controle sobre documentos, consultas e dependências externas.

Este trabalho também se justifica pela comparação planejada com grandes modelos de linguagem. Modelos como ChatGPT, Gemini e Claude são capazes de responder a muitas perguntas gerais, mas, sem acesso ao manual correto, podem generalizar informações. A hipótese deste trabalho é que um sistema RAG, mesmo utilizando um modelo local menor, pode apresentar vantagem em perguntas dependentes de documentação específica.

\section{Metodologia}
\label{sec:metodologia}

A metodologia adotada combina revisão bibliográfica, desenvolvimento de artefato computacional e planejamento de avaliação experimental. A revisão abrange modelos de linguagem, busca semântica, RAG, FAISS, recuperação multimodal, CLIP, memória conversacional e prevenção de alucinações. A etapa prática consistiu na implementação de uma aplicação composta por frontend, backend, scripts Python, índice vetorial, modelo local e persistência de conversas.

O desenvolvimento ocorreu de forma incremental. Inicialmente, a recuperação era baseada principalmente em busca vetorial. Em seguida, foram adicionados sinais lexicais, regras de domínio, metadados de veículo, memória conversacional e validação visual. As mudanças mais recentes buscaram tornar o fluxo mais coeso: o sistema passou a reescrever perguntas incompletas, rejeitar imagens não validadas, aplicar pós-processamento em respostas frágeis e construir respostas determinísticas em casos de localização de componentes quando há evidência clara no manual.

A avaliação comparativa com grandes modelos de linguagem será conduzida após a estabilização do sistema. A proposta é construir um conjunto de perguntas automotivas com respostas de referência baseadas nos manuais, executar o sistema RAG local e comparar suas respostas com as de modelos sem acesso aos documentos. A comparação deverá considerar correção factual, utilidade, citação de fonte, completude e presença de alucinações.

\section{Organização do trabalho}

O Capítulo~\ref{cap:revisao} apresenta a fundamentação teórica e os trabalhos relacionados. O Capítulo~\ref{cap:desenvolvimento} descreve a arquitetura e a implementação da solução. O Capítulo~\ref{cap:resultados} discute os testes realizados, o estágio atual da validação e o desenho da avaliação comparativa. Por fim, o Capítulo~\ref{cap:conclusao} apresenta as considerações finais e os próximos passos do projeto.
